\section{Reproducibility and Artifact Statement}

The repository pins Python dependencies, the base-model revision, static data,
source commits, report hashes, and the generated-evaluation paths. The local
training stack was Python 3.12.3, PyTorch 2.13.0, torchvision 0.28.0,
Transformers 5.14.1, TRL 1.9.2, PEFT 0.20.0, Datasets 5.0.1, Accelerate 1.14.0,
and Trackio 0.34.0. Runs used CUDA 13.0 and cuDNN 92000 on an NVIDIA GeForce RTX
5070 Laptop GPU (compute capability 12.0) with 8,546,484,224 bytes of reported
VRAM. A fixed seed and fixed greedy decoding reduce avoidable variation, but
PyTorch does not guarantee bit-identical execution across releases, platforms,
or devices. \citep{transformersgeneration,pytorchreproducibility}
The framework and experiment-tracking citations identify the corresponding
software families; the exact installed versions above come from the evaluation
records. \citep{paszke2019pytorch,wolf2020transformers,vonwerra2020trl,
mangrulkar2022peft,trackio}
\claimsource{eval-collection}\claimsource{manifest}

The public pipeline enforced a GitHub-first gate before baseline generation or
training: clean synchronized main, tracked-source presence, ignored/untracked
credential file, and a scan proving the credential bytes did not occur in Git
objects. Generated reports were allowlist-validated and paired JSON/Markdown
was rendered from one structured object. Failing checkpoints remained ignored
operational state. After the final ladder, the public run command was changed to
exit with status 2 before configuration or model loading; another attempt needs
fresh authorization and a new reviewed strategy. \citep{gitcatfile}
The upload interface existed behind the acceptance gate but was not invoked.
\citep{huggingfacehub}
\claimsource{code-gitgate}\claimsource{code-reporting}
\claimsource{source-failclosed}

Source changes and evidence landed in separate public merge-history stages:
foundation PR~\href{https://github.com/BurnyCoder/training-facts-into-llms/pull/1}{\#1},
paper adaptation PR~\href{https://github.com/BurnyCoder/training-facts-into-llms/pull/2}{\#2},
initial evidence PRs~\href{https://github.com/BurnyCoder/training-facts-into-llms/pull/3}{\#3}--
\href{https://github.com/BurnyCoder/training-facts-into-llms/pull/4}{\#4},
semantic source/evidence PRs~\href{https://github.com/BurnyCoder/training-facts-into-llms/pull/5}{\#5}--
\href{https://github.com/BurnyCoder/training-facts-into-llms/pull/6}{\#6}, and
minimal-pair source/evidence PRs~\href{https://github.com/BurnyCoder/training-facts-into-llms/pull/7}{\#7}--
\href{https://github.com/BurnyCoder/training-facts-into-llms/pull/8}{\#8}.
Appendix~\ref{app:run-evidence} gives immutable commit and evidence links.
The repository was originally named \texttt{fact-teaching}; GitHub now
canonicalizes it as \texttt{BurnyCoder/training-facts-into-llms}.
\claimsource{pr-foundation}\claimsource{pr-paper}\claimsource{pr-semantic}
\claimsource{pr-minimal}\claimsource{pr-results}

The machine-readable source of truth is
\texttt{reports/manifest.json}\claimsource{manifest}. The complete technical
retrospective is
\texttt{reports/EXPERIMENTS.md}\claimsource{retrospective}. The LaTeX
source and named PDF are derived artifacts: their CPU contract checks that all
nine run bindings, eight results, interruption, evidence paths, and negative
publication state remain synchronized. The paper can be rebuilt with
\code{make -C paper}; build intermediates are ignored and the stable PDF is
copied to \code{output/pdf/teaching-one-synthetic-fact-qwen35.pdf}.
\citep{uv}\claimsource{paper-build}\claimsource{paper-contract}
